% Montage de distillation (exercice 10, chapitre 9)
\begin{tikzpicture}[x=1cm,y=1cm]
% Chauffe-ballon et ballon
\filldraw[pcAppareil,rounded corners=1.6mm] (-1.15,-0.1) rectangle (1.15,0.85);
\fill[yellow!14] (-0.72,0.9) arc[start angle=180,end angle=360,radius=0.72] -- cycle;
\draw[pcVerre] (-0.72,0.9) arc[start angle=180,end angle=360,radius=0.72];
\draw[pcVerre] (-0.72,0.9) .. controls (-0.5,1.35) .. (-0.2,1.52) (0.72,0.9) .. controls (0.5,1.35) .. (0.2,1.52);
% Colonne et tête de distillation
\draw[pcVerre] (-0.2,1.52) -- (-0.2,3.6) (0.2,1.52) -- (0.2,3.15);
\draw[pcVerre] (0.2,3.15) -- (0.75,2.95) (0.2,3.6) -- (0.2,3.42) -- (0.85,3.18);
\draw[pcVerre] (-0.2,3.6) -- (0.2,3.6);
% Thermomètre
\draw[otInk!70,line width=0.8pt] (0,3.4) -- (0,4.4);
\fill[otRed!80] (0,3.3) circle (0.06);
% Réfrigérant droit incliné
\begin{scope}[shift={(0.8,3.06)},rotate=-22]
  \filldraw[fill=pcEau,draw=otNavy,line width=0.6pt] (0.3,-0.2) rectangle (2.9,0.2);
  \draw[pcVerre] (0,-0.07) -- (3.3,-0.07) (0,0.07) -- (3.3,0.07);
  \draw[pcFil] (0.55,0.2) -- (0.55,0.6) (2.65,-0.2) -- (2.65,-0.6);
  \coordinate (sortie) at (3.3,0); \coordinate (entree) at (2.65,-0.6); \coordinate (haut) at (0.55,0.6);
\end{scope}
\draw[pcFlecheN] ($(entree)+(0.6,0)$) -- (entree);
\draw[pcFlecheN] (haut) -- ++(-0.5,0.15);
% Allonge et erlenmeyer collecteur
\draw[pcVerre] (sortie) -- ++(0.15,-0.55);
\pic at ($(sortie)+(0.25,-2.75)$) {erlen={hauteur=1.9,niveau=0.35,liquide=yellow!25}};
\pcNote[east]{(0,4.2)}{(-0.9,4.35)}{thermomètre}
\pcNote[west]{(2.1,2.75)}{(2.6,3.4)}{réfrigérant droit}
\node[pcLegende,anchor=west] at ($(entree)+(0.65,0)$) {entrée d'eau};
\pcNote[east]{(-0.7,1.0)}{(-1.4,1.4)}{ballon : mélange}
\pcNote[west]{($(sortie)+(0.55,-2.5)$)}{($(sortie)+(1.2,-2.2)$)}{distillat}
\end{tikzpicture}
\par\smallskip\legende{L'espèce la plus volatile se vaporise la première ; ses vapeurs se condensent dans le réfrigérant et sont recueillies.}
